\section*{Lectures on Lie Groups and Lie Algebras – 2,3}

\section*{More Examples of Lie Algebras}

The matrix algebra, considered as a Lie algebra with respect to the commutator,
is denoted $\mathfrak{gl}_n(\mathbb{C})$ (or simply $\mathfrak{gl}_n$).

The following subspaces of $\mathfrak{gl}_n$ are Lie subalgebras (i.e., closed under the commutator):

\begin{enumerate}
\item The subspace of scalar matrices $\mathbb{C}E$ (this is moreover an ideal in $\mathfrak{gl}_n$),
\item The subspace of traceless matrices $\mathfrak{sl}_n$ (also an ideal),
\item The subspace of strictly upper triangular matrices $\mathfrak{n}_+(n)$,
\item The subspace of upper triangular matrices $\mathfrak{b}_+(n)$,
\item The subspace of skew-symmetric matrices $\mathfrak{so}_n$.
\end{enumerate}

\textbf{Proposition 1.}
$\mathfrak{n}_+(3)$ is the Heisenberg algebra, and
$\mathfrak{b}_+(2)$ is the direct sum of a one-dimensional and a two-dimensional non-abelian Lie algebra.

\section*{Representations: Glossary}

\textbf{Definition 1.}
A \emph{representation} of a Lie algebra $L$ in a vector space $V$ is a homomorphism
\[
\rho : L \to \mathfrak{gl}(V).
\]

\textbf{Definition 2.}
A representation $(V,\rho)$ is called \emph{faithful} if $\rho$ is injective.

\textbf{Definition 3.}
A \emph{subrepresentation} of $(V,\rho)$ is a subspace $V' \subset V$ invariant under all operators $\rho(x)$.
A representation is \emph{irreducible} if its only subrepresentations are $0$ and $V$.

\textbf{Definition 4.}
If $(V_1,\rho_1)$ and $(V_2,\rho_2)$ are representations,
their direct sum is defined by
\[
\rho(x)(v_1 \oplus v_2) = \rho_1(x)v_1 \oplus \rho_2(x)v_2.
\]
A representation is \emph{indecomposable} if it cannot be written as a direct sum of proper subrepresentations.

\textbf{Definition 5.}
A representation is \emph{completely reducible} if every subrepresentation has an invariant complement.

\textbf{Proposition 2.}
A finite-dimensional completely reducible representation decomposes as a direct sum of irreducibles.

\section*{Examples of Representations}

\begin{enumerate}
\item The trivial representation: $\rho(x)=0$.

\item The adjoint representation: $\rho(x)=\mathrm{ad}\,x$.

Subrepresentations in the adjoint representation are precisely ideals.
$L$ is simple iff the adjoint representation is irreducible.
It is faithful iff the center is trivial.

\item The tautological representation of $\mathfrak{gl}_n$ on $V=\mathbb{C}^n$.
It is always faithful.
\end{enumerate}

\textbf{Proposition 5.}
The tautological representation of $\mathfrak{gl}_n$ is irreducible.

\textbf{Proposition 6.}
The tautological representation of $\mathfrak{n}_+(2)$ is reducible but indecomposable.

\section*{Invariant Forms}

\textbf{Definition 6.}
A bilinear form $\Phi$ on $(V,\rho)$ is \emph{invariant} if
\[
\Phi(\rho(x)v_1,v_2)+\Phi(v_1,\rho(x)v_2)=0.
\]

Examples:

\begin{enumerate}
\item Standard inner product on $\mathfrak{so}_n$.
\item On $\mathfrak{gl}_n$, $\Phi(A,B)=\mathrm{tr}(AB)$.
\item Killing form:
\[
K(x,y)=\mathrm{tr}(\mathrm{ad}\,x\,\mathrm{ad}\,y).
\]
\end{enumerate}

Orthogonal complements of subrepresentations are subrepresentations.

If the invariant form is positive definite (real case), the representation is completely reducible.

\section*{Homomorphisms of Representations}

\textbf{Definition 7.}
$\varphi : V_1 \to V_2$ is an intertwining operator if
\[
\varphi \rho_1(x) = \rho_2(x) \varphi.
\]

\textbf{Proposition 8 (Schur's Lemma).}
If $\varphi$ is a homomorphism between irreducible representations,
then either $\varphi=0$ or $\varphi$ is an isomorphism.

Over $\mathbb{C}$, any endomorphism of an irreducible representation is scalar multiplication.

\section*{Tensor Products and Duals}

\textbf{Definition 8.}
The tensor product representation is defined by
\[
\rho(x)(v_1 \otimes v_2)
=
\rho_1(x)v_1 \otimes v_2
+
v_1 \otimes \rho_2(x)v_2.
\]

\textbf{Proposition 11.}
$S^n V$ and $\Lambda^n V$ are subrepresentations of $V^{\otimes n}$.

\textbf{Definition 9.}
The dual representation is defined by
\[
\rho^*(x)f(v) = -f(\rho(x)v).
\]

\section*{Canonical Isomorphisms}

\textbf{Proposition 13.}
For $V=\mathbb{C}^n$,
\[
\text{Adjoint representation of } \mathfrak{gl}_n
\cong V \otimes V^*.
\]

\textbf{Definition 10.}
The invariant subspace:
\[
V^L = \{ v \in V \mid \rho(x)v=0 \ \forall x\in L \}.
\]

\[
\mathrm{Hom}_L(V,W) = (V^* \otimes W)^L
\]

Invariant bilinear forms correspond to $(V^* \otimes V^*)^L$.

\section*{Irreducible Representations of $\mathfrak{sl}_2$}

\textbf{Proposition 16.}
The symmetric powers $S^n V$ of the tautological representation of $\mathfrak{sl}_2$
are irreducible.

Let $V=\mathbb{C}^2$ with basis $\{u,v\}$ and

\[
e =
\begin{pmatrix}
0 & 1 \\
0 & 0
\end{pmatrix},
\quad
f =
\begin{pmatrix}
0 & 0 \\
1 & 0
\end{pmatrix},
\quad
h =
\begin{pmatrix}
1 & 0 \\
0 & -1
\end{pmatrix}.
\]

Then on $S^n V$ (homogeneous polynomials of degree $n$):

\[
\rho(e) = u \frac{\partial}{\partial v},
\quad
\rho(f) = v \frac{\partial}{\partial u},
\quad
\rho(h) = u \frac{\partial}{\partial u}
-
v \frac{\partial}{\partial v}.
\]

Irreducibility follows from:

\begin{enumerate}
\item Every nontrivial invariant subspace contains $u^n$.
\item $u^n$ generates $S^n V$.
\end{enumerate}